% Canonical Paper I source.
The regular-AES axioms admit a homogeneous realization on the positive affine
group.  This construction explains the upper-half-plane metric from the group
action and the normalization of the arithmetic frame; it does not assert that
every regular AES is homogeneous or that the model is unique.

\subsection{The positive affine group}

Write
\[
 G=\operatorname{Aff}^{+}(1,\mathbb R)
   =\{(x,y):x\in\mathbb R,\ y>0\}
\]
with multiplication
\begin{equation}\label{eq:positive-affine-group-law}
 (x,y)(x',y')=(x+yx',yy').
\end{equation}
The element \((x,y)\) represents the affine map \(t\mapsto x+yt\).
Left translation by \((x_0,y_0)\) is
\[
 L_{(x_0,y_0)}(x,y)=(x_0+y_0x,y_0y).
\]
For fixed \(\mu,\lambda>0\), consider the left-invariant vector fields
\begin{equation}\label{eq:hyperbolic-arithmetic-frame}
 X_u=-\mu y\,\partial_x,
 \qquad
 X_v=-\lambda y\,\partial_y .
\end{equation}
Both signs are chosen so that the frame is positively oriented for
\(dx\wedge dy\) and has the desired action on the assignment below.

\begin{proposition}[Invariant affine metric]\label{prop:invariant-affine-metric}
The Riemannian metric for which the frame
\eqref{eq:hyperbolic-arithmetic-frame} is orthonormal is
\begin{equation}\label{eq:hyperbolic-aes-metric}
 g_{\mu,\lambda}
 =\frac1{y^2}
  \left(\frac{dx^2}{\mu^2}+\frac{dy^2}{\lambda^2}\right).
\end{equation}
It is invariant under every left translation of \(G\).
\end{proposition}

\begin{proof}
The coframe dual to \((X_u,X_v)\) is
\[
 -\frac{dx}{\mu y},
 \qquad
 -\frac{dy}{\lambda y},
\]
whose sum of squares is \eqref{eq:hyperbolic-aes-metric}.  Under
\(L_{(x_0,y_0)}\), one has
\[
 x\longmapsto x_0+y_0x,\qquad
 y\longmapsto y_0y,\qquad
 dx\longmapsto y_0dx,\qquad
 dy\longmapsto y_0dy.
\]
Substitution in \eqref{eq:hyperbolic-aes-metric} shows
\(L_{(x_0,y_0)}^*g_{\mu,\lambda}=g_{\mu,\lambda}\).
\end{proof}

\subsection{The basic hyperbolic arithmetic expression space}

Let
\[
 \mathbb H^2=\{(x,y)\in\mathbb R^2:y>0\}
\]
with its standard orientation, and define
\begin{equation}\label{eq:basic-hyperbolic-assignment}
 a(x,y)=-\frac{x}{y}.
\end{equation}

\begin{theorem}[Basic hyperbolic AES]\label{thm:basic-hyperbolic-aes}
For every \(\mu,\lambda>0\), the tuple
\[
 \mathfrak E_{\mathrm{hyp}}(\mu,\lambda)
 =(\mathbb H^2,g_{\mu,\lambda},a;\mu,\lambda)
\]
is a regular arithmetic expression space.  Its canonical arithmetic frame is
the left-invariant frame \eqref{eq:hyperbolic-arithmetic-frame}.  Moreover,
\((\mathbb H^2,g_{\mu,\lambda})\) is complete and has constant Gaussian
curvature
\[
 K=-\lambda^2.
\]
\end{theorem}

\begin{proof}
Direct differentiation gives
\[
 a_x=-\frac1y,
 \qquad
 a_y=\frac{x}{y^2}=-\frac ay.
\]
Therefore the frame \eqref{eq:hyperbolic-arithmetic-frame} satisfies
\[
 X_u a=\mu,
 \qquad
 X_v a=\lambda a.
\]
It is positively oriented and orthonormal by
Proposition~\ref{prop:invariant-affine-metric}.  The frame equivalence of
Proposition~\ref{prop:canonical-arithmetic-frame} now proves the regular-AES
identity.  Equivalently, using
\[
 g_{\mu,\lambda}^{-1}
 =\mu^2y^2\,\partial_x\otimes\partial_x
  +\lambda^2y^2\,\partial_y\otimes\partial_y,
\]
one obtains directly
\[
 |\nabla a|_{g_{\mu,\lambda}}^2
 =\mu^2y^2a_x^2+\lambda^2y^2a_y^2
 =\mu^2+\lambda^2a^2.
\]

For completeness and curvature, introduce the global coordinate
\[
 X=\frac{\lambda}{\mu}x.
\]
Then
\begin{equation}\label{eq:scaled-standard-hyperbolic-metric}
 g_{\mu,\lambda}
 =\frac1{\lambda^2}\frac{dX^2+dy^2}{y^2}.
\end{equation}
The metric in parentheses is the complete standard upper-half-plane metric of
curvature \(-1\).  Multiplication of a Riemannian metric by
\(\lambda^{-2}\) preserves completeness and multiplies Gaussian curvature by
\(\lambda^2\).  Hence \(g_{\mu,\lambda}\) is complete and has curvature
\(-\lambda^2\).
\end{proof}

The parameters have different geometric roles.  The change
\(X=(\lambda/\mu)x\) removes \(\mu\) from the metric, so \(\mu\) fixes the
additive normalization of the assignment frame.  The parameter \(\lambda\)
sets the curvature scale.  No uniqueness claim for regular AESs follows from
this normalization.

\subsection{Horocyclic coordinates and assignment contours}

Set \(y=e^{\lambda v}\) and \(u=x\).  Then
\begin{equation}\label{eq:horocyclic-affine-coordinates}
 g_{\mu,\lambda}
 =e^{-2\lambda v}\frac{du^2}{\mu^2}+dv^2,
 \qquad
 a(u,v)=-u e^{-\lambda v}.
\end{equation}
The curves \(y=\mathrm{constant}\) are horocycles based at the ideal point
\(\infty\), and the vertical lines \(x=\mathrm{constant}\) are their
orthogonal geodesics.  The assignment contour of value \(c\) is
\[
 a^{-1}(c)=\{x=-cy,\ y>0\}.
\]
In particular, the zero locus in this model is the vertical geodesic
\(\{x=0,\ y>0\}\).

\subsection{The Laplace eigenfunction}

We use the divergence-of-gradient sign convention
\begin{equation}\label{eq:laplacian-sign-convention}
 \Delta_g f
 =\frac1{\sqrt{\det g}}\,
  \partial_i\!\left(\sqrt{\det g}\,g^{ij}\partial_jf\right).
\end{equation}

\begin{proposition}[Laplace eigenfunction]\label{prop:laplace-eigenfunction}
The assignment function of the basic hyperbolic AES satisfies
\begin{equation}\label{eq:hyperbolic-assignment-laplacian}
 \Delta_{g_{\mu,\lambda}}a=2\lambda^2a.
\end{equation}
\end{proposition}

\begin{proof}
Since
\[
 \sqrt{\det g_{\mu,\lambda}}=\frac1{\mu\lambda y^2},
\]
the products
\(\sqrt{\det g}\,g^{xx}=\mu/\lambda\) and
\(\sqrt{\det g}\,g^{yy}=\lambda/\mu\) are constant.  Thus
\[
 \Delta_{g_{\mu,\lambda}}f
 =\mu^2y^2 f_{xx}+\lambda^2y^2 f_{yy}.
\]
For \(a=-x/y\), one has \(a_{xx}=0\) and
\(a_{yy}=-2x/y^3=2a/y^2\).  Substitution gives
\eqref{eq:hyperbolic-assignment-laplacian}.
\end{proof}

\subsection{Assignment-compatible grid transformations}

For \(s\in\mathbb R\) and \(k>0\), define
\begin{equation}\label{eq:assignment-grid-transformations}
 \mathsf X_s(x,y)=(x-sy,y),
 \qquad
 \mathsf Y_k(x,y)=\left(x,\frac yk\right).
\end{equation}
They implement addition and positive multiplication at the level of the
assignment:
\begin{equation}\label{eq:grid-assignment-action}
 a\circ\mathsf X_s=a+s,
 \qquad
 a\circ\mathsf Y_k=ka.
\end{equation}
Figure~\ref{fig:hyperbolic-arithmetic-grid} separates the assignment contours
from the horocyclic background and also marks the two assignment-compatible
motions.  The pullback calculation following the figure determines their
metric status.

\begin{figure}[t]
\centering
\begin{tikzpicture}[
  x=1.15cm,y=1.05cm,
  axis/.style={-{Latex[length=2mm]},thick,black!70},
  gridline/.style={thin,gray!35},
  contour/.style={very thick,red!65!black},
  move/.style={-{Latex[length=2.2mm]},very thick,blue!65!black},
  every node/.style={font=\small}
]
  \fill[blue!2] (-3.25,0) rectangle (3.25,3.15);
  \draw[axis] (-3.4,0) -- (3.5,0) node[right] {$x$};
  \draw[axis] (0,-0.05) -- (0,3.35) node[above] {$y$};
  \node[anchor=north west,font=\scriptsize] at (-3.25,-0.08) {$y=0$};

  \foreach \yy in {0.65,1.3,1.95,2.6}
    {\draw[gridline] (-3.2,\yy) -- (3.2,\yy);}
  \foreach \xx in {-2.4,-1.2,1.2,2.4}
    {\draw[gridline] (\xx,0.05) -- (\xx,3.1);}
  \draw[contour] (0,0.03) -- (-3.05,3.08);
  \draw[contour] (0,0.03) -- (0,3.1);
  \draw[contour] (0,0.03) -- (3.05,3.08);
  \node[red!65!black,anchor=south east,font=\scriptsize] at (-3.0,3.05)
    {$a=1$};
  \node[red!65!black,anchor=south west,font=\scriptsize] at (0.05,3.0)
    {$a=0$};
  \node[red!65!black,anchor=south west,font=\scriptsize] at (3.0,3.05)
    {$a=-1$};
  \node[red!65!black,font=\scriptsize,fill=white,inner sep=1pt]
    at (-1.75,1.15) {contour $a=c$: $x=-cy$};

  \coordinate (p) at (1.85,2.3);
  \coordinate (px) at (0.75,2.3);
  \coordinate (py) at (1.85,1.15);
  \fill (p) circle (1.6pt) node[above right] {$p$};
  \draw[move] (p) -- (px)
    node[midway,above,fill=white,inner sep=1pt,font=\scriptsize]
    {$\mathsf X_s:\ a\mapsto a+s$};
  \draw[move,dashed] (p) -- (py)
    node[midway,right,fill=white,inner sep=1pt,font=\scriptsize]
    {$\mathsf Y_k:\ a\mapsto ka$};
  \fill[blue!65!black] (px) circle (1.3pt);
  \fill[blue!65!black] (py) circle (1.3pt);
\end{tikzpicture}
\caption{The basic hyperbolic arithmetic grid on the upper half-plane.  Gray
lines show the horocycle--geodesic background; red rays are the level sets
$a=-x/y$.  The blue motions implement addition and positive multiplication on
the assignment.  They preserve the indicated assignment rules but are not, in
general, isometries of $g_{\mu,\lambda}$.}
\label{fig:hyperbolic-arithmetic-grid}
\end{figure}

These maps are \emph{assignment-compatible transformations}; in general they
are not isometries of \(g_{\mu,\lambda}\).  Indeed,
\begin{align}
 \mathsf X_s^*g_{\mu,\lambda}
 &=\frac1{y^2}
   \left(\frac{(dx-s\,dy)^2}{\mu^2}
         +\frac{dy^2}{\lambda^2}\right),\label{eq:additive-grid-pullback}\\
 \mathsf Y_k^*g_{\mu,\lambda}
 &=\frac1{y^2}
   \left(\frac{k^2dx^2}{\mu^2}
         +\frac{dy^2}{\lambda^2}\right).
 \label{eq:multiplicative-grid-pullback}
\end{align}
Consequently \(\mathsf X_s\) is an isometry only for \(s=0\), and
\(\mathsf Y_k\) only for \(k=1\).

\begin{remark}[A Baumslag--Solitar relation]\label{rem:hyperbolic-bs-relation}
For an integer \(n\ge2\),
\(\mathsf X_s^n=\mathsf X_{ns}\), and direct composition gives
\begin{equation}\label{eq:hyperbolic-bs-relation}
 \mathsf Y_n^{-1}\circ\mathsf X_s^n\circ\mathsf Y_n
 =\mathsf X_s.
\end{equation}
With stable letter \(T=\mathsf Y_n^{-1}\), this reads
\(T\mathsf X_s^nT^{-1}=\mathsf X_s\), a
Baumslag--Solitar relation in the \(BS(n,1)\) convention.  Equation
\eqref{eq:hyperbolic-bs-relation} is an identity among
assignment-compatible diffeomorphisms, not a claim that the arithmetic grid acts
by hyperbolic isometries.  No complexity conclusion is inferred from it.
\end{remark}

Coordinate derivations of the group chart, gradient, curvature, Laplacian, metric
pullbacks, and relation \eqref{eq:hyperbolic-bs-relation} are collected in
Appendix~\ref{app:hyperbolic-calculations}.
